\section*{Chapter \ref{ch:g8:project-data-explorer} --- Project: A Data Explorer}

\begin{solution}{exo:g8:project-data-explorer:1}
For the 5, it prints \code{please type a number from 0 to 4}; for the 1,
it prints \code{question 1 is not ready yet}; for the 0, \code{goodbye}. A
test of the chapter plays this session.
\end{solution}

\begin{solution}{exo:g8:project-data-explorer:2}
Birch and Cedar have 2 loans each, Ash 1. The keys are \code{[-2, "Birch"]},
\code{[-2, "Cedar"]} and \code{[-1, "Ash"]}: Birch, then Cedar (same first
\omterm{def:g7:lists:list}{item}, \code{"Birch"} before \code{"Cedar"}), then Ash.
\end{solution}

\begin{solution}{exo:g8:project-data-explorer:3}
\omcode{code/grade-8/06-project-data-explorer/exercises/by_author.py}{1}{7}{The
books of one author.}\label{lst:g8:project-data-explorer:sol-author}
\end{solution}

\begin{solution}{exo:g8:project-data-explorer:4}
The rows are built from the loans, one row per loan: a book borrowed 10 times
gets 10 identical rows, all at the top. The rows must be built from the books,
one per book, as in version 4.
\end{solution}

\begin{solution}{exo:g8:project-data-explorer:5}
\omcode{code/grade-8/06-project-data-explorer/exercises/busiest_month.py}{1}{11}{The
busiest month.}\label{lst:g8:project-data-explorer:sol-month}

On the town's table it returns \code{('09', 45)}: September, with 45 loans.
\end{solution}

\begin{solution}{exo:g8:project-data-explorer:6}
Every date of the table starts with \code{2}, and \code{"30/6/2025"} with
\code{3}. As text, every out date is therefore smaller: every loan seems to
have gone out before that ``date''. No return date is larger, so only the
loans that never came back are counted: the 22 of
\cref{ch:g8:files-and-data-tables}. The \omterm{def:g4:programs-and-operating-system:program}{program} gives a plausible number, and
it is wrong; that is why the date is checked first.
\end{solution}

\begin{solution}{exo:g8:project-data-explorer:7}
The same pattern as \code{never_borrowed}, on members:

\omcode{code/grade-8/06-project-data-explorer/exercises/never_borrowers.py}{1}{10}{Members
without a loan.}\label{lst:g8:project-data-explorer:sol-members}

On the town's tables it returns an empty \omterm{def:g7:lists:list}{list}: all 60 members borrowed at
least once.
\end{solution}

\begin{solution}{exo:g8:project-data-explorer:8}
\code{book_finder} writes 80 times, once per book. Then, for each of the 400
loans, \code{loans_by_genre} reads \code{finder} once, reads \code{counts}
once with \code{get}, and writes it once: $3 \times 400 = 1{,}200$. In all,
$80 + 1{,}200 = 1{,}280$ steps, each a single step whatever the size of the
\omterm{def:g8:dictionaries:dictionary}{dictionary}. Searching the \omterm{def:g7:lists:list}{list} of books for each loan instead could take up
to 80 \omterm{def:g6:if-else:comparison}{comparisons} per loan.
\end{solution}

\begin{solution}{exo:g8:project-data-explorer:9}
\code{csv.DictWriter} takes the \omterm{def:g8:files-and-data-tables:table}{header row} from the keys of the first row:

\omcode{code/grade-8/06-project-data-explorer/exercises/save_csv.py}{1}{10}{The
top ten as a \omterm{def:g8:files-and-data-tables:table}{CSV file}.}\label{lst:g8:project-data-explorer:sol-csv}
\end{solution}

\begin{solution}{exo:g8:project-data-explorer:10}
\omcode{code/grade-8/06-project-data-explorer/exercises/loan_length.py}{1}{14}{The
average length of a loan.}\label{lst:g8:project-data-explorer:sol-length}

On the 378 loans that came back, the average is 17.6 days.
\end{solution}

\begin{solution}{exo:g8:project-data-explorer:11}
Each loan adds 1 to the count of its book, wherever it is in the \omterm{def:g3:files-and-folders:file}{file}, so the
counts do not depend on the order. The rows are built in the order of the
books, and then \omterm{def:g8:searching-and-sorting:sorted}{sorted} by the key \code{[-loans, title]}. The 80 titles are all
different, so no two rows have the same key: the \omterm{def:g8:searching-and-sorting:sorted}{sorted} order is decided by
the keys alone, and the first ten rows are always the same.
\end{solution}

\begin{solution}{exo:g8:project-data-explorer:12}
First the \omterm{def:g7:lists:list}{list} of different books of each member; then, for each member,
every pair of those books, counted in a \omterm{def:g8:dictionaries:dictionary}{dictionary} whose keys are texts such
as \code{"5+47"}:

\omcode{code/grade-8/06-project-data-explorer/exercises/best_pair.py}{1}{21}{The
pair borrowed by the most members.}\label{lst:g8:project-data-explorer:sol-pair}

It returns \code{('5+47', 3)}: three members borrowed both books 5 and 47.
Other pairs also reach 3; the \omterm{def:g7:functions:function}{function} keeps the first one it meets. Sorting
the numbers of each member's books with \code{key=int} makes sure each pair is
always written the same way, smaller number first.
\end{solution}
